\documentclass[11pt]{article}
\usepackage{amsmath,amssymb}
\usepackage{booktabs}
\usepackage{graphicx}
\usepackage[margin=2.2cm]{geometry}

\title{Test de simbiosis \LaTeX{} + Python en Calc}
\author{}
\date{\today}

\begin{document}
\maketitle

Cada bloque \verb|%#python ... %#end| se ejecuta en el kernel (no se imprime);
los resultados entran en el documento con \verb|\py{...}|.

%#python
import math
import numpy as np
import pandas as pd
import sympy as sp
import matplotlib
import matplotlib.pyplot as plt

# --- Aritmética básica ---
g, t = 9.81, 2.5
v = g * t
d = 0.5 * g * t**2
fact10 = math.factorial(10)
print(f"v={v:.2f}  d={d:.2f}  10!={fact10}")
%#end

\section{Aritmética y \texttt{math}}
Caída libre: $v = g t = \py{round(v,2)}\ \mathrm{m/s}$ y
$d = \tfrac12 g t^2 = \py{round(d,2)}\ \mathrm{m}$.
Además $10! = \py{fact10}$, $\pi \approx \py{round(math.pi,6)}$,
$e \approx \py{round(math.e,6)}$ y una operación compuesta
$v\cdot 2 - 3/2 = \py{round(v*2 - 3/2, 3)}$.

%#python
# --- NumPy: álgebra lineal ---
A = np.array([[2, 1, -1], [-3, -1, 2], [-2, 1, 2]], dtype=float)
b = np.array([8, -11, -3], dtype=float)
x = np.linalg.solve(A, b)
detA = np.linalg.det(A)
eig = np.linalg.eigvals(A)

def mat_latex(M, fmt="{:.0f}"):
    filas = [" & ".join(fmt.format(c) for c in fila) for fila in M]
    return r"\begin{pmatrix}" + r" \\ ".join(filas) + r"\end{pmatrix}"

A_tex = mat_latex(A)
x_tex = ",\\ ".join(f"{xi:.3f}" for xi in x)
eig_tex = ",\\ ".join(f"{e.real:.2f}" for e in eig)
%#end

\section{NumPy --- álgebra lineal}
Resolviendo $A\mathbf{x}=\mathbf{b}$ con
\[ A = \py{A_tex}, \qquad \det(A) = \py{round(detA,1)}. \]
La solución es $\mathbf{x} = (\py{x_tex})$ y los autovalores de $A$ son
$\{\py{eig_tex}\}$.

%#python
# --- pandas: tabla a LaTeX ---
df = pd.DataFrame({
    "$n$":        list(range(1, 6)),
    "$n^2$":      [n**2 for n in range(1, 6)],
    "$\\sqrt{n}$": [round(math.sqrt(n), 4) for n in range(1, 6)],
    "$\\ln n$":   [round(math.log(n), 4) for n in range(1, 6)],
})
tabla_tex = df.to_latex(index=False, escape=False, column_format="cccc")
%#end

\section{Tabla generada con pandas}
\begin{center}
\py{tabla_tex}
\end{center}

%#python
# --- SymPy: cálculo simbólico ---
xs = sp.symbols('x')
f = sp.sin(xs) * sp.exp(-xs)
fprime = sp.simplify(sp.diff(f, xs))
F = sp.integrate(f, xs)
raices = sp.solve(xs**2 - 5*xs + 6, xs)

f_tex = sp.latex(f)
fp_tex = sp.latex(fprime)
F_tex = sp.latex(F)
raices_tex = ",\\ ".join(sp.latex(r) for r in raices)
%#end

\section{SymPy --- cálculo simbólico}
Sea $f(x) = \py{f_tex}$. Entonces
\[ f'(x) = \py{fp_tex}, \qquad \int f(x)\,dx = \py{F_tex} + C. \]
Las raíces de $x^2 - 5x + 6 = 0$ son $x \in \{\py{raices_tex}\}$.

%#python
# --- matplotlib: figuras (se redimensionan con \includegraphics) ---
fig, ax = plt.subplots(figsize=(6, 3.4))
xx = np.linspace(0, 4*np.pi, 400)
ax.plot(xx, np.sin(xx), label="sin(x)")
ax.plot(xx, np.cos(xx), "--", label="cos(x)")
ax.plot(xx, np.exp(-xx/8)*np.sin(xx), ":", label="amortiguada")
ax.legend(); ax.grid(True, alpha=0.3)
ax.set_title("Funciones")
ruta_graf = figure("trig", dpi=200)
plt.close(fig)

fig2, ax2 = plt.subplots(figsize=(4.5, 3))
ax2.bar(["A", "B", "C", "D"], [23, 45, 12, 38],
        color=["#4c72b0", "#dd8452", "#55a868", "#c44e52"])
ax2.set_title("Barras")
ruta_barras = figure("barras", dpi=200)
plt.close(fig2)
%#end

\section{Gráficas con matplotlib (imágenes redimensionables)}
La misma figura al 80\% del ancho de texto:
\begin{center}\includegraphics[width=0.8\textwidth]{\py{ruta_graf}}\end{center}

Dos imágenes con tamaños distintos, una rotada $90^\circ$:
\begin{center}
\includegraphics[width=0.30\textwidth]{\py{ruta_barras}}\qquad
\includegraphics[height=3cm, angle=90]{\py{ruta_barras}}
\end{center}

%#python
# --- Estadística (numpy) ---
datos = [4.2, 5.1, 3.8, 6.0, 5.5, 4.9, 5.3, 6.1, 4.4, 5.0]
media = float(np.mean(datos))
desv = float(np.std(datos, ddof=1))
mediana = float(np.median(datos))
datos_tex = ",\\ ".join(str(v) for v in datos)
%#end

\section{Estadística}
Para los $n = \py{len(datos)}$ datos $\{\py{datos_tex}\}$:
$\bar{x} = \py{round(media,3)}$, $s = \py{round(desv,3)}$,
mediana $= \py{round(mediana,2)}$.

\end{document}
